% ============================================================
% Chapter 10 --- The chain rule and directional derivatives
% Originally: content/week-04.tex, \section{Chain rule} with its subsections
%             \subsection{Directional derivatives} and
%             \subsection{What does imply differentiability} (Corsin Week 4, Monday),
%             plus \section{The chain rule along a path} (Corsin Week 4, Friday).
%             The two subsections are promoted to sections here.
% ============================================================
\chapter{The Chain Rule and Directional Derivatives}
\label{ch:chain_rule}

% Transition: Claude Opus 5 (High)
Having encoded the differential as a matrix, we now ask how it interacts with composition. The
answer is the chain rule, and in this setting it says something cleaner than its one-variable
version does: the differential of a composition is the \emph{composition of the differentials},
which is to say the product of their matrices. Reading it that way makes it a statement of linear
algebra rather than a formula to be memorised, and explains where the sums of products in the
coordinate form come from.

The same machinery yields directional derivatives, and with them the standard trap of
several-variable calculus. A function can possess every directional derivative at a point and
still fail to be differentiable there, or even continuous.
\Cref{sec:sufficient_condition_differentiability} says what does suffice, and it is the criterion
used in practice for the rest of the course.

% Originally: content/week-04.tex, \section{Chain rule} (Corsin Week 4, Monday)

% Source: Corsin Nick/Class Notes/Week 4.pdf, p. 5
\section{Chain rule}
\label{sec:chain_rule}

\begin{theorem}[Chain rule]
\label{thm:chain_rule}
The \newterm{chain rule} \germanfor{Kettenregel} states: for
$f : U \subseteq \mathbb{R}^n \to \mathbb{R}^m$, $g : V \subseteq \mathbb{R}^m \to \mathbb{R}^d$
differentiable, it holds:
\[D(g\circ f)(x) = Dg(f(x))\,Df(x).\]
The same identity holds for the corresponding Jacobi matrices, where the composition on the
right becomes an ordinary matrix product.
\end{theorem}

% Generator: Claude Opus 5 (High)
\begin{ainote}
Corsin states the chain rule without proof, as does the lecture at this point. The one-line
reason it is true is worth carrying, though: differentiability says $f$ and $g$ are each well
approximated by a linear map near the relevant point, and composing two approximations composes
the linear maps, with the error terms of the two compositions still $o(|h|)$. Making that
precise is the whole proof, and it is the same argument as in one variable.
\end{ainote}

% Supplement: Sascha Brack/Class Notes/Week_04_Notes_Friday_Updated.pdf, pp. 7--10
\begin{remark}[The chain rule as a sum over paths]
\label{rem:chain_rule_paths}
Stated as $D(g\circ f) = Dg\cdot Df$ the rule is compact but, as Sascha Brack puts it in his
notes, \qt{pretty abstract}. It says nothing about \emph{why} a sum appears when the rule is
written out in coordinates. His device is to draw the variables as a graph and read the formula
off it.

Write $x \in \mathbb{R}^n \xmapsto{\ f\ } y \in \mathbb{R}^m \xmapsto{\ g\ } z \in \mathbb{R}^k$
and join each variable to those it feeds into. Then $x_i$ influences $z_j$ only by first moving
some intermediate $y_\ell$, and the contribution of that route is the product of the two rates
along it. Summing over the routes gives the chain rule:

% Sonnet 5 (Medium) --- FIG-W04-03, after the variable graph on Sascha Brack's p. 10.
\begin{center}
\begin{tikzpicture}[>=Latex, scale=1.0]
  \node (xi) at (0,2.4) {$x_i$};
  \node (y1) at (-2.2,1.1) {$y_1$};
  \node (y2) at (-0.7,1.1) {$y_2$};
  \node (yd) at (0.75,1.1) {$\cdots$};
  \node (ym) at (2.2,1.1) {$y_m$};
  \node (zj) at (0,-0.2) {$z_j$};

  \foreach \t in {y1,y2,ym}{\draw[->, blue!70!black] (xi) -- (\t);}
  \foreach \s in {y1,y2,ym}{\draw[->, orange!90!black] (\s) -- (zj);}

  \node[blue!70!black, font=\scriptsize, above left] at (-1.25,1.75)
    {$\dfrac{\partial y_\ell}{\partial x_i}$};
  \node[orange!90!black, font=\scriptsize, below left] at (-1.25,0.55)
    {$\dfrac{\partial z_j}{\partial y_\ell}$};
\end{tikzpicture}
\end{center}

\[\frac{\partial z_j}{\partial x_i}
  \;=\; \sum_{\ell=1}^{m}
  \underbrace{\frac{\partial z_j}{\partial y_\ell}}_{(\Jac g)_{j\ell}}
  \underbrace{\frac{\partial y_\ell}{\partial x_i}}_{(\Jac f)_{\ell i}}
  \;=\; \big(\Jac g(f(x))\,\Jac f(x)\big)_{ji}.\]

The right-hand equality is just the definition of matrix multiplication, and that is precisely the
point: \textbf{matrix multiplication is a sum over paths}, and the chain rule looks like a product of
matrices precisely because composing functions composes those graphs. Sascha's summary is that
computing $\partial z_j/\partial x_i$ this way is nothing more than \qt{computing single
components of the chain rule}.

The reading is worth having because it makes the coordinate form reconstructible instead of
memorised: count the intermediate variables, one term each, each term a product of two factors,
and the index that is summed over is the one belonging to the middle layer. It is also the
form in which the rule is actually used in physics. Compare
\Cref{sec:chain_rule_along_path}, which is this picture with a single intermediate layer
$\gamma_1,\dots,\gamma_n$ and a single source variable $t$.
\end{remark}

% Generator: Claude Opus 5 (High) --- not in any source read so far; added because it is the
% cleanest one-line payoff of the chain rule along a path and is quoted constantly in physics.
\begin{theorem}[Euler's homogeneous function theorem]
\label{thm:eulers_homogeneous_function_theorem}
Let $f \in C^1(\mathbb{R}^n\setminus\{0\},\mathbb{R})$ be \newterm{positively homogeneous of
degree $k$} for some $k \in \mathbb{R}$, meaning that
\[f(tx) = t^kf(x) \qquad \text{for all } x \in \mathbb{R}^n\setminus\{0\} \text{ and all } t>0.\]
Then $\langle x, \nabla f(x)\rangle = k\,f(x)$ for all $x \in \mathbb{R}^n\setminus\{0\}$.
\end{theorem}

\begin{proof}
Fix $x \neq 0$ and read both sides of the homogeneity identity as functions of $t>0$ alone. The
left-hand side is $f$ composed with the path $\gamma(t) := tx$, whose velocity is
$\gamma'(t) = x$, so the chain rule along a path (\cref{sec:chain_rule_along_path}) gives
\[\frac{d}{dt}f(tx) = \big\langle \nabla f(tx),\ x\big\rangle.\]
On the right-hand side $f(x)$ is a constant in $t$, and therefore
$\frac{d}{dt}\big(t^kf(x)\big) = kt^{k-1}f(x)$. The two derivatives agree for every $t>0$;
evaluating at $t=1$ leaves $\langle \nabla f(x), x\rangle = k f(x)$, which is the claim by
symmetry of the inner product.
\end{proof}

\begin{aiexample}[Chain Rule for Vector Fields]
% Generator: Gemini 3.6 Flash (Medium)
Let $f(u,v) := u^2 + v^2$ and let $g(r,\theta) := (r\cos\theta, r\sin\theta)\transp$. The composite function is $(f \circ g)(r,\theta) = r^2\cos^2\theta + r^2\sin^2\theta = r^2$.
By the multi-variable Chain Rule, the Jacobian matrix of $f \circ g$ satisfies:
\[
D(f \circ g)(r,\theta) = Df(g(r,\theta)) \cdot Dg(r,\theta) = \begin{pmatrix} 2r\cos\theta & 2r\sin\theta \end{pmatrix} \begin{pmatrix} \cos\theta & -r\sin\theta \\ \sin\theta & r\cos\theta \end{pmatrix} = \begin{pmatrix} 2r & 0 \end{pmatrix},
\]
which matches the direct calculation $\frac{\partial}{\partial r}(r^2) = 2r$ and $\frac{\partial}{\partial \theta}(r^2) = 0$.
\end{aiexample}

% Source: Corsin Nick/Class Notes/Week 4.pdf, pp. 5--6
\begin{exercise}[Jacobian of a composite map]
\label{ex:chain_rule_computation}
Compute $\Jac(g\circ f)$ for
\[f : \mathbb{R}^3\to\mathbb{R}^2,\ x \mapsto (x_1x_2,\ x_2+x_3), \qquad g : \mathbb{R}^2\to\mathbb{R}^2,\ x \mapsto (e^{x_1},\ x_1x_2).\]
\end{exercise}

% Kept inline: solved directly in Corsin Nick/Class Notes/Week 4.pdf, pp. 5--6.
\begin{exercisesolution}[\cref{ex:chain_rule_computation}]
We have
\[\Jac f(x) = \begin{pmatrix} x_2 & x_1 & 0 \\ 0 & 1 & 1\end{pmatrix}, \qquad \Jac g(x) = \begin{pmatrix} e^{x_1} & 0 \\ x_2 & x_1 \end{pmatrix}.\]
So:
\[
\begin{aligned}
\Jac(g\circ f) &= \begin{pmatrix} e^{x_1x_2} & 0 \\ x_2+x_3 & x_1x_2\end{pmatrix}\begin{pmatrix} x_2 & x_1 & 0 \\ 0 & 1 & 1\end{pmatrix} \\
&= \begin{pmatrix} x_2e^{x_1x_2} & x_1e^{x_1x_2} & 0 \\ x_2^2 + x_2x_3 & 2x_1x_2 + x_1x_3 & x_1x_2 \end{pmatrix}
\end{aligned}
\]
\end{exercisesolution}

% Source: Jérôme Paschoud/Notizen/Woche 4.2 Kettenregel.pdf, p. 1
\begin{exercise}[Computing differentials with and without the chain rule]
\label{ex:chain_rule_computation_paschoud}
Let $f: \mathbb{R}^2 \to \mathbb{R}^3$ and $g: \mathbb{R}^3 \to \mathbb{R}$ be defined by
\[f(x,y) := \begin{pmatrix} x+y \\ x-y \\ x^2+y^2-1 \end{pmatrix}, \qquad g(x,y,z) := x^2+y^2+z^2.\]
Compute the differential $D(g \circ f)$ in two ways:
\begin{enumerate}[label=\textbf{(\alph*)}]
  \item By directly computing the composition $g \circ f$ and then differentiating.
  \item By computing $Df$ and $Dg$ and using the chain rule.
\end{enumerate}
\end{exercise}

\begin{exercisesolution}[\cref{ex:chain_rule_computation_paschoud}]
\begin{enumerate}[label=\textbf{(\alph*)}]
  \item \textbf{Directly:} First compute $g(f(x,y))$:
    \begin{align*}
    (g \circ f)(x,y) &= (x+y)^2 + (x-y)^2 + (x^2+y^2-1)^2 \\
    &= (x^2+2xy+y^2) + (x^2-2xy+y^2) + (x^4+y^4+1+2x^2y^2-2x^2-2y^2) \\
    &= 2x^2 + 2y^2 + x^4 + y^4 + 1 + 2x^2y^2 - 2x^2 - 2y^2 \\
    &= x^4 + y^4 + 2x^2y^2 + 1 \\
    &= (x^2+y^2)^2 + 1.
    \end{align*}
    Now differentiate this scalar function with respect to $x$ and $y$:
    \[D(g \circ f)(x,y) = \begin{pmatrix} 2(x^2+y^2)\cdot 2x & 2(x^2+y^2)\cdot 2y \end{pmatrix} = \begin{pmatrix} 4x(x^2+y^2) & 4y(x^2+y^2) \end{pmatrix}.\]

  \item \textbf{Via chain rule:} Compute the differentials of $f$ and $g$ individually.
    \[Df(x,y) = \begin{pmatrix} 1 & 1 \\ 1 & -1 \\ 2x & 2y \end{pmatrix}, \qquad Dg(u,v,w) = \begin{pmatrix} 2u & 2v & 2w \end{pmatrix}.\]
    Evaluate $Dg$ at $f(x,y)$:
    \[Dg(f(x,y)) = \begin{pmatrix} 2(x+y) & 2(x-y) & 2(x^2+y^2-1) \end{pmatrix}.\]
    Multiply the matrices:
    \begin{align*}
    D(g \circ f)(x,y) &= Dg(f(x,y)) Df(x,y) \\
    &= \begin{pmatrix} 2(x+y) & 2(x-y) & 2(x^2+y^2-1) \end{pmatrix} \begin{pmatrix} 1 & 1 \\ 1 & -1 \\ 2x & 2y \end{pmatrix} \\
    &= \begin{pmatrix} 4x + 4x(x^2+y^2-1) & 4y + 4y(x^2+y^2-1) \end{pmatrix} \\
    &= \begin{pmatrix} 4x(x^2+y^2) & 4y(x^2+y^2) \end{pmatrix}.
    \end{align*}
    Both methods yield the same result.
\end{enumerate}
\end{exercisesolution}

% Source: old_exams/examHS24.pdf (13 February 2025), Wahr/Falsch-Fragen, Aufgabe 3, p. 2
% Extractor: Gemini 3.1 Pro (High)
\begin{exercise}[True or False \textnormal{(the higher-order chain rule)}]
\label{ex:examHS24_TF_chain_rule}
Let $u \in C^2(\mathbb{R}^n)$ and $F = (F_1, \dots, F_n) \in C^2(\mathbb{R}^n, \mathbb{R}^n)$, with
$n \ge 2$. Decide whether each of the following identities is true or false.
\begin{enumerate}[label=\textbf{(\alph*)}]
  \item The identity $\partial_i(u \circ F) = \sum_{j=1}^n \partial_ju \circ F \, \partial_iF_j$ holds for all $i \in \{1, \dots, n\}$.
  \item The identity $\partial_{ii}(u \circ F) = \sum_{j=1}^n \partial_{jj}u \circ F \, (\partial_iF_j)^2 + \sum_{j=1}^n \partial_ju \circ F \, \partial_{ii}F_j$ holds for all $i \in \{1, \dots, n\}$.
  \item The identity $\partial_i\partial_j(u \circ F) = \partial_j\partial_i(u \circ F)$ holds for all $i, j \in \{1, \dots, n\}$.
\end{enumerate}
\exinfo{This exercise is taken from the exam of 13 February 2025 (Prof.\ Joaquim Serra), where it appears as the third
block of true/false questions.}
\exsol{sol:chain_rule_higher_order_tf}
\end{exercise}

% Source: old_exams/examFS24.pdf (21 August 2024), Wahr/Falsch-Fragen, Aufgabe 3, pp. 2--3
% Extractor: Claude Opus 5 (High)
% Presented as a worked example rather than as an exercise: the block is three one-line verdicts,
% and its value is entirely in seeing which hypothesis each one consumes. Its companion above,
% ex:examHS24_TF_chain_rule, is the same theme set as a problem, and assumes F in C^2 where this
% one assumes only C^1 --- which is the whole of part (c).
\begin{example}[Which identities are forced, and by what]
\label{ex:forced_identities_chain_rule}
% Generator: Claude Opus 5 (High) --- the three statements are the examiner's; the discussion is
% written for these notes.
Let $u \in C^2(\mathbb{R}^n)$ and let $F = (F_1,\dots,F_n) \in C^1(\mathbb{R}^n,\mathbb{R}^n)$,
with $n \geq 2$. Exactly one of the following three identities is forced by those hypotheses, and
the two that fail do so for quite different reasons.
\begin{enumerate}[label=\textbf{(\alph*)}]
  \item \textbf{$\partial_iF_j = \partial_jF_i$ for all $i,j$: not forced.} This asks the Jacobi
    matrix $\Jac F$ to be symmetric at every point, which is the integrability condition of
    \cref{lem:integrability_conditions}: it holds precisely when $F$ is a candidate for a gradient,
    and nothing in the hypotheses makes it one. The rotation field
    $F(x) := (-x_2,\, x_1,\, 0,\dots,0)$ is as smooth as one likes and has $\partial_1F_2 = 1$
    while $\partial_2F_1 = -1$.
  \item \textbf{$\partial_j\partial_iu = \partial_i\partial_ju$ for all $i,j$: forced.} This is
    Schwarz's lemma (\cref{thm:schwarz_mixed_partials}), and $u \in C^2$ is exactly its hypothesis.
    Note how differently the two statements read once they are set side by side: \textbf{(a)}
    constrains an arbitrary vector field, \textbf{(b)} constrains the gradient of a function. Only
    the second is automatic.
  \item \textbf{$u \circ F \in C^2$: not forced.} The chain rule gives
    $\partial_i(u\circ F) = \sum_{j}(\partial_ju\circ F)\,\partial_iF_j$, and the factors
    $\partial_iF_j$ are merely continuous, so this says only that $u \circ F \in C^1$. To
    differentiate once more one would have to differentiate $\partial_iF_j$, which the hypotheses
    do not permit. Take $n = 2$, $u(y_1,y_2) := y_1$ and $F(x_1,x_2) := (x_1\lvert x_1\rvert,\, x_2)$.
    Then $F \in C^1$, since $\partial_1F_1 = 2\lvert x_1\rvert$ is continuous, and
    $(u\circ F)(x) = x_1\lvert x_1\rvert$ has the same non-differentiable first derivative at
    $x_1 = 0$. Consequently $u \circ F \notin C^2$.
\end{enumerate}
The general rule that \textbf{(c)} illustrates is that a composition inherits the \emph{weaker} of
the two regularities. Raise $F$ to $C^2$ and $u \circ F$ becomes $C^2$, which is what
\cref{ex:examHS24_TF_chain_rule}\textbf{(c)} is then free to assume.
\exinfo{This example is taken from the exam of 21 August 2024 (Prof.\ Joaquim Serra), where it appears as the third block
of true/false questions.}
\end{example}

% Originally: content/week-04.tex, \subsection{Directional derivatives} --- promoted to a section

% Sonnet 5 (Medium)
\section{Directional derivatives}
\label{sec:directional_derivatives}

% Source: Corsin Nick/Class Notes/Week 4.pdf, p. 6
Sometimes the \newterm{differential} (\germanterm{das Differential}) is not easily expressed as a
Jacobian. Then it can be convenient to use the formula for \newterm{directional derivatives}
(\germanterm{Richtungsableitung}):
\begin{equation*}
DF_{x_0}(v) = \partial_v F(x_0) = \left.\frac{d}{ds}\right|_{s=0} F(x_0 + sv) \tag{$*$}
\end{equation*}

% Generator: Claude Opus 5 (Medium) --- the geometry behind (*): pick a ray, restrict to it,
% and you are back to a one-variable derivative from Analysis I.
Formula $(*)$ describes a two-step recipe, and the picture below is worth carrying around.
\textbf{Step one:} pick a direction $v$ and walk away from $x_0$ along the straight ray
$s \mapsto x_0+sv$, ignoring the rest of the domain. \textbf{Step two:} record the values of $F$
along that ray. What you get is an ordinary function of \emph{one} variable $s$, and
$\partial_vF(x_0)$ is simply its derivative at $s=0$ --- an Analysis I object. The partial
derivative $\partial_iF$ is the special case $v = e_i$: walk parallel to the $i$-th axis.

\begin{center}
\begin{tikzpicture}[>=Latex, scale=1.0]
  % --- Panel 1: the domain, and the rays we are allowed to walk along ---
  \begin{scope}
    \draw[->] (-0.3,0) -- (2.9,0) node[right, font=\small] {$x_1$};
    \draw[->] (0,-0.3) -- (0,2.7) node[above, font=\small] {$x_2$};

    \coordinate (P) at (1.15,1.0);
    % the rays, extended in grey
    \draw[gray!55, dashed] (-0.15,1.0) -- (2.75,1.0);
    \draw[gray!55, dashed] (0.15,0.15) -- (2.35,2.35);

    \draw[orange!90!black, very thick, ->] (P) -- ++(0.8,0) node[below, font=\scriptsize] {$e_1$};
    \draw[green!55!black, very thick, ->] (P) -- ++(0,0.8) node[left, font=\scriptsize] {$e_2$};
    \draw[purple, very thick, ->] (P) -- ++(0.62,0.62) node[above left, font=\scriptsize] {$v$};

    \fill (P) circle (1.8pt) node[below left, font=\scriptsize] {$x_0$};
    \node[below, font=\small, align=center] at (1.3,-0.75)
      {pick a direction and\\ walk along that ray};
  \end{scope}

  % --- Panel 2: the profile along each ray, and its slope at s=0 ---
  \begin{scope}[shift={(6.6,0)}]
    \draw[->] (-1.1,0) -- (2.2,0) node[right, font=\small] {$s$};
    \draw[->] (0,-0.3) -- (0,2.8) node[above right, font=\small] {$F(x_0+sv)$};
    \draw (0,0.05) -- (0,-0.05) node[below right, font=\scriptsize] {$0$};

    % profile along e_1 (gentle positive slope) and along v (negative slope)
    \draw[orange!90!black, very thick, domain=-0.55:1.95, samples=50, smooth]
      plot ({\x}, {1.4 + 0.35*\x + 0.18*\x*\x});
    \draw[purple, very thick, domain=-0.55:1.55, samples=50, smooth]
      plot ({\x}, {1.4 - 0.75*\x + 0.30*\x*\x});

    % their tangents at s = 0
    \draw[orange!90!black, dashed] (-0.5,1.225) -- (1.95,2.0825);
    \draw[purple, dashed] (-0.5,1.775) -- (1.55,0.2375);

    \fill (0,1.4) circle (1.8pt);
    \draw[dotted, gray] (-0.62,1.4) -- (0,1.4);
    \node[font=\scriptsize, left] at (-0.65,1.4) {$F(x_0)$};

    % anchored just off the right end of each tangent line, clear of both curves
    \node[orange!90!black, font=\scriptsize, right] at (2.0,2.20) {slope $=\partial_1F(x_0)$};
    \node[purple, font=\scriptsize, right] at (1.62,0.22) {slope $=\partial_vF(x_0)$};

    \node[below, font=\small, align=center] at (0.5,-0.75)
      {each ray gives a one-variable function;\\ its slope at $0$ is $\partial_vF(x_0)$};
  \end{scope}
\end{tikzpicture}
\end{center}

The right-hand panel is also the cleanest way to see what differentiability adds. Every direction
produces \emph{some} slope; nothing so far forces those slopes to be related to one another. Being
differentiable is precisely the demand that all of them are read off a \emph{single} linear map
via $\partial_vF(x_0) = DF_{x_0}(v)$ --- so that knowing the $n$ partial derivatives already
determines every other direction. \Cref{ex:all_directional_derivatives_exist} below is the
cautionary case where all the slopes exist but no such linear map governs them.

\begin{importantremark}
\label{rem:directional_derivative_caveat}
The directional derivative $\partial_v F(x_0)$ can exist (depending on $v$ and $x_0$) even if $F$
is not differentiable at $x_0$, i.e.\ $DF_{x_0}$ does not exist. However, if $DF_{x_0}$ exists,
the formula $(*)$ holds.
\end{importantremark}

% Generator: Claude Opus 5 (Medium)
\begin{aiexample}[Every direction, and still not continuous]
\label{ex:all_directional_derivatives_exist}
\Cref{ex:partial_derivatives_not_continuous} showed that the two \emph{partial} derivatives are
too little. One might hope that asking for \emph{all} directions is enough. It is not. Let
\[f(x,y) := \begin{cases} \dfrac{x^2y}{x^4+y^2} & (x,y) \neq (0,0), \\[1ex] 0 & (x,y)=(0,0).\end{cases}\]
Fix a direction $v = (a,b) \neq 0$ and compute along the ray $t \mapsto tv$:
\[\frac{f(ta,tb)}{t} = \frac{1}{t}\cdot\frac{t^2a^2\cdot tb}{t^4a^4+t^2b^2}
   = \frac{a^2b}{t^2a^4+b^2} \;\xrightarrow{\ t\to0\ }\;
   \begin{cases} \dfrac{a^2}{b} & b \neq 0, \\[1ex] 0 & b = 0.\end{cases}\]
So $\partial_vf(0,0)$ exists for \textbf{every} direction $v$. Yet $f$ is not even continuous at
the origin: approaching along the \emph{parabola} $y=x^2$ instead of along a straight line,
\[f(x,x^2) = \frac{x^2\cdot x^2}{x^4+x^4} = \frac12 \quad\text{for all } x \neq 0.\]
By \cref{prop:differentiable_implies_continuous}, $f$ is therefore not differentiable at the
origin.

Two lessons worth keeping.
\begin{itemize}
  \item \textbf{Straight lines cannot see everything.} Every directional derivative is computed
    along a line, and this $f$ is well behaved along every line while misbehaving along a
    parabola. Approaching a point in $\mathbb{R}^n$ is genuinely harder than in $\mathbb{R}$,
    where there are only two directions.
  \item \textbf{Watch whether $v \mapsto \partial_vf(x_0)$ is linear.} If $f$ were differentiable
    we would need $\partial_vf(0,0) = DF_0(v)$, a \emph{linear} function of $v$; but $a^2/b$ is
    not linear in $(a,b)$. That failure alone rules out differentiability, without any mention of
    continuity --- and it is precisely the test that settles \cref{ex:5.1}\textbf{(b)} on the
    \cref{ex:5.3}.
\end{itemize}
\end{aiexample}

% Originally: content/week-04.tex, \subsection{What does imply differentiability} --- promoted

% Supplement: Sascha Brack/Class Notes/Week_04_Notes_Monday_Updated.pdf, pp. 3--5
\section{What \emph{does} imply differentiability}
\label{sec:sufficient_condition_differentiability}

\begin{ainote}
Corsin's notes give the definition and the warnings, but never state a \emph{sufficient}
condition --- so far we have two counterexamples and no positive test. Sascha Brack's notes
(Lecture 8) state one as Theorem 2.9 and organise the whole picture into a chain of
implications; both are reproduced below, since without them the practical question
\qt{how do I actually check differentiability?} is left unanswered.
\end{ainote}

\begin{theorem}[Continuous partial derivatives suffice]
\label{thm:continuous_partials_imply_differentiable}
Let $U \subseteq \mathbb{R}^n$ be open and $F : U \to \mathbb{R}^m$. If all partial derivatives
$\partial_iF$ exist on $U$ \textbf{and are continuous} at every $x \in U$, then $F$ is
differentiable at every $x \in U$, and
\[DF_x = \big(\partial_1F(x),\ \dots,\ \partial_nF(x)\big).\]
Such an $F$ is called \newterm{continuously differentiable}, written $F \in C^1(U)$.
\end{theorem}

This is the theorem that makes the subject workable: in practice one almost never checks the
definition directly. Compute the partials, observe that they are built from continuous pieces,
and differentiability follows --- which is exactly why the course can write $f \in C^1$, $C^2$,
$C^k$ so freely from here onward.

Together with the two counterexamples, the whole landscape is a single chain, each link
one-directional:

% Sonnet 5 (Medium) --- FIG-W04-02, after the diagram on Sascha Brack's p. 3.
\begin{center}
\begin{tikzpicture}[>=Latex, node distance=0pt]
  \tikzset{
    box/.style={draw, rounded corners=2pt, align=center, font=\small,
                inner sep=4pt, minimum height=8mm},
    imp/.style={->, very thick, green!45!black},
    nimp/.style={->, thick, red!75!black, dashed}}

  \node[box, fill=green!8] (c1) at (0,4.5) {$F \in C^1$ near $x_0$\\[-2pt]
    {\scriptsize (all $\partial_iF$ exist \emph{and} are continuous)}};
  \node[box, fill=blue!8] (df) at (0,2.9) {$F$ differentiable at $x_0$};
  \node[box] (dir) at (0,1.4) {all directional derivatives $\partial_vF(x_0)$ exist};
  \node[box] (par) at (0,-0.1) {all partial derivatives $\partial_iF(x_0)$ exist};

  \draw[imp] (c1) -- (df);
  \draw[imp] (df) -- (dir);
  \draw[imp] (dir) -- (par);

  \node[font=\scriptsize, green!45!black, right] at (0.15,3.7)
    {\cref{thm:continuous_partials_imply_differentiable}};
  \node[font=\scriptsize, green!45!black, right] at (0.15,2.15) {formula $(*)$};
  \node[font=\scriptsize, green!45!black, right] at (0.15,0.65) {take $v = e_i$};

  % The reverse arrows, all false. Each carries its counterexample as a midway node on the
  % arc itself, filled white so it interrupts the dashed line cleanly. Previously the three
  % labels were free-standing nodes anchored at x = -2.6 while the arcs bulged out past
  % x = -3, so every arrow was drawn straight through its own label.
  \tikzset{noLbl/.style={midway, font=\scriptsize, red!75!black, align=right,
                         fill=white, inner sep=1.5pt}}
  \draw[nimp] (df.west)  to[out=160, in=200, looseness=1.7]
    node[noLbl] {\textbf{no:} \cref{ex:differentiable_not_c1}} (c1.west);
  \draw[nimp] (dir.west) to[out=160, in=200, looseness=1.7]
    node[noLbl] {\textbf{no:} \cref{ex:all_directional_derivatives_exist}} (df.west);
  \draw[nimp] (par.west) to[out=160, in=200, looseness=1.7]
    node[noLbl] {\textbf{no:} \cref{ex:partial_derivatives_not_continuous}} (dir.west);
\end{tikzpicture}
\end{center}

Every downward arrow holds, and \textbf{not one of them reverses} --- each failure is witnessed
by an example already in this chapter, so the diagram doubles as a map of where those
counterexamples belong:
\begin{itemize}
  \item \textbf{Differentiable $\not\Rightarrow$ $C^1$:} \cref{ex:differentiable_not_c1}, whose
    partial derivatives oscillate without limit at the origin even though the differential
    exists there.
  \item \textbf{All directional derivatives $\not\Rightarrow$ differentiable:}
    \cref{ex:all_directional_derivatives_exist}, where every ray gives a slope but
    $v \mapsto \partial_vf(0,0)$ is not linear.
  \item \textbf{Partial derivatives $\not\Rightarrow$ all directional derivatives:}
    \cref{ex:partial_derivatives_not_continuous} again. For $f = \tfrac{xy}{x^2+y^2}$ and a
    direction $v = (a,b)$ with $ab \neq 0$,
    \[\frac{f(ta,tb)}{t} = \frac{1}{t}\cdot\frac{t^2ab}{t^2(a^2+b^2)}
      = \frac{ab}{t\,(a^2+b^2)} \ \xrightarrow{\ t \to 0\ }\ \pm\infty,\]
    so $\partial_vf(0,0)$ exists only along the two axes --- the two directions the partial
    derivatives happen to look at.
\end{itemize}
The three counterexamples are not variations on one theme; they fail at three different
heights, which is precisely why the chain has three separate links.

% Supplement: Sascha Brack/Class Notes/Week_04_Notes_Monday_Updated.pdf, p. 5
\begin{remark}[A recipe for checking differentiability at a point]
\label{rem:differentiability_recipe}
Sascha's notes distil the above into three steps, worth following in this order.
\begin{enumerate}[label=\textbf{\arabic*.}]
  \item \textbf{Compute the partial derivatives $\partial_iF(x_0)$.} If one of them fails to
    exist, $F$ is not differentiable and you are done. If they all exist \emph{and are
    continuous near $x_0$}, $F$ \emph{is} differentiable by
    \Cref{thm:continuous_partials_imply_differentiable} --- also done. Most exercises end here.
  \item \textbf{Otherwise, guess the candidate.} If $F$ is differentiable at all, the differential
    can only be $L := \big(\partial_1F(x_0),\dots,\partial_nF(x_0)\big)$, since the partials are
    forced (\cref{sec:jacobi_matrix}).
  \item \textbf{Verify the candidate against the definition:}
    \[\lim_{h\to0}\frac{\big|F(x_0+h)-F(x_0)-L(h)\big|}{\lVert h\rVert} \overset{?}{=} 0.\]
    In $\mathbb{R}^2$, substituting $h = (r\cos\theta,\ r\sin\theta)$ and checking that the
    result tends to $0$ \emph{uniformly in $\theta$} is usually the quickest route --- the same
    polar trick used in \cref{ex:all_directional_derivatives_exist} and, later, in
    \Cref{q:quiz_q4}.
\end{enumerate}
Step 1 is the one students skip; it settles the overwhelming majority of cases without any
$\varepsilon$--$\delta$ work at all.
\end{remark}

% Supplement: Sascha Brack/Class Notes/Week_03_Notes_Friday.pdf, p. 7
\begin{example}[Checking differentiability at the origin]
\label{ex:differentiability_recipe_application}
Let $f(x,y) := \frac{x^2 y}{x^2+y^2}$ for $(x,y) \neq (0,0)$ and $f(0,0) := 0$. We check if $f$ is differentiable at the origin.
\begin{enumerate}[label=\textbf{\arabic*.}]
  \item \textbf{Partial derivatives:}
    \[\partial_1 f(0,0) = \lim_{t\to 0}\frac{f(t,0)-f(0,0)}{t} = \lim_{t\to 0}\frac{0}{t} = 0, \qquad \partial_2 f(0,0) = \lim_{t\to 0}\frac{f(0,t)-f(0,0)}{t} = 0.\]
    They exist, but to use \cref{thm:continuous_partials_imply_differentiable} we would have to compute them everywhere and show they are continuous at $(0,0)$, which is tedious and (as we will see) false. We proceed to step 2.
  \item \textbf{Guess the candidate:} The only possible differential is $L(h_1,h_2) = \partial_1 f(0,0)h_1 + \partial_2 f(0,0)h_2 = 0$.
  \item \textbf{Verify against the definition:} We check whether the following limit is zero:
    \[\lim_{h\to 0} \frac{|f(h_1,h_2) - f(0,0) - L(h)|}{\lVert h\rVert} = \lim_{h\to 0} \frac{\frac{h_1^2 h_2}{h_1^2+h_2^2}}{\sqrt{h_1^2+h_2^2}} = \lim_{h\to 0} \frac{h_1^2 h_2}{(h_1^2+h_2^2)^{3/2}}.\]
    Passing to polar coordinates $h = (r\cos\theta, r\sin\theta)$, the expression becomes
    \[\frac{r^3\cos^2\theta\sin\theta}{r^3} = \cos^2\theta\sin\theta.\]
    This depends on the angle of approach $\theta$ (for instance, it is $0$ along the axes but $\frac{1}{2\sqrt{2}} \neq 0$ along $h_1=h_2$), so the limit as $r \to 0$ is not uniformly $0$. Hence $f$ is \textbf{not differentiable} at $(0,0)$.
\end{enumerate}
\end{example}

% Supplement: Sascha Brack/Class Notes/Week_04_Notes_Friday_Updated.pdf, pp. 4--5
\begin{example}[A function that is differentiable at the origin]
\label{ex:differentiability_recipe_positive}
Let $f(x,y) := \frac{x^2 y^2}{\sqrt{x^2+y^2}}$ for $(x,y) \neq (0,0)$ and $f(0,0) := 0$. We check differentiability at the origin following the recipe.
\begin{enumerate}[label=\textbf{\arabic*.}]
  \item \textbf{Partial derivatives:} For $(x,y) \neq (0,0)$, the quotient and chain rules give
    \[\partial_1 f(x,y) = xy^2\frac{x^2+2y^2}{(x^2+y^2)^{3/2}}, \qquad \partial_2 f(x,y) = x^2y\frac{2x^2+y^2}{(x^2+y^2)^{3/2}}.\]
    At the origin, computing directly from the definition yields $\partial_1 f(0,0) = 0$ and $\partial_2 f(0,0) = 0$.
    To see if they are continuous at $(0,0)$, we substitute $x = r\cos\theta$ and $y = r\sin\theta$:
    \[\partial_1 f(r\cos\theta, r\sin\theta) = \frac{r\cos\theta\,r^2\sin^2\theta \cdot r^2(1+\sin^2\theta)}{r^3} = r^2\cos\theta\sin^2\theta(1+\sin^2\theta).\]
    As $r \to 0$, this tends to $0$ uniformly in $\theta$, so $\partial_1 f$ is continuous at $(0,0)$. The same holds for $\partial_2 f$. By \cref{thm:continuous_partials_imply_differentiable}, $f$ is differentiable at the origin.
  \item \textbf{Guess the candidate:} We already know $f$ is differentiable and $L = (0,0)$. But for practice, let's verify it with the definition directly.
  \item \textbf{Verify against the definition:} We check the limit
    \[\lim_{h\to 0} \frac{|f(h_1,h_2) - f(0,0) - L(h)|}{\lVert h\rVert} = \lim_{h\to 0} \frac{\frac{h_1^2 h_2^2}{\sqrt{h_1^2+h_2^2}}}{\sqrt{h_1^2+h_2^2}} = \lim_{h\to 0} \frac{h_1^2 h_2^2}{h_1^2+h_2^2}.\]
    In polar coordinates, this is $\frac{r^4\cos^2\theta\sin^2\theta}{r^2} = r^2\cos^2\theta\sin^2\theta$. As $r \to 0$, the limit is clearly $0$ uniformly in $\theta$, confirming $f$ is differentiable with $Df_0 = 0$.
\end{enumerate}
\end{example}

% Originally: content/week-04.tex, \section{The chain rule along a path} (Corsin Week 4, Friday)

% Source: Corsin Nick/Class Notes/Week 4.pdf, p. 8
\section{The chain rule along a path}
\label{sec:chain_rule_along_path}

An important special case of the chain rule is that of a composition of a scalar field with a
path. Suppose $f : \mathbb{R}^n \supseteq U \to \mathbb{R}$ and
$\gamma : [0,1]\to\mathbb{R}^n$ are a $C^1$ \newterm{scalar field} (\germanterm{Skalarfeld}) and
path, respectively. Then
\[
\begin{aligned}
D(f\circ\gamma)(t) &= Df(\gamma(t))\cdot D\gamma(t) \\
&= \langle \nabla f(\gamma(t)),\ \gamma'(t)\rangle \\
&= \sum_{i=1}^{n} \frac{\partial f}{\partial \gamma_i}\frac{\partial \gamma_i}{\partial t}
\end{aligned}
\]
where $\gamma = (\gamma_1,\dots,\gamma_n)$ and
$\dfrac{\partial f}{\partial \gamma_i} := \partial_i f(\gamma(t))$.

This operation is so common in physics that it has its own notation: for $f$ as above, we look at
each coordinate as a function of \textbf{time} and differentiate:
\[\frac{d}{dt}f(x_1,\dots,x_n) := Df(x_1(t),\dots,x_n(t)) = \sum_{i=1}^{n}\frac{\partial f}{\partial x_i}\frac{\partial x_i}{\partial t}.\]

% Generator: Claude Opus 5 (High)
The middle line, $D(f\circ\gamma)(t) = \langle\nabla f(\gamma(t)),\gamma'(t)\rangle$, is worth
dwelling on, because it is the formula that gives the gradient its meaning. It says the rate at
which $f$ changes along a path depends on the velocity only through its inner product with
$\nabla f$ --- so by Cauchy--Schwarz (\cref{thm:cauchy_schwarz}) the change is fastest when
$\gamma'$ points along $\nabla f$, and zero when $\gamma'$ is perpendicular to it. Two standard
facts fall straight out:

\begin{itemize}
  \item \textbf{$\nabla f$ points in the direction of steepest ascent}, with
    $|\nabla f|$ the rate of ascent in that direction.
  \item \textbf{$\nabla f$ is perpendicular to the level sets of $f$.} If $\gamma$ stays inside
    a level set $\{f = c\}$, then $f\circ\gamma$ is constant, so
    $\langle\nabla f(\gamma(t)),\gamma'(t)\rangle = 0$ for every $t$ --- and $\gamma'$ runs
    through the directions tangent to the level set.
\end{itemize}

\begin{center}
% Generator: Claude Opus 5 (High) --- FIG-W04-04: level sets of f(x,y) = x^2 + 2y^2, a path
% crossing them, and the decomposition of gamma' at one point.
% Level sets are the ellipses x^2 + 2y^2 = c, i.e. semi-axes sqrt(c) and sqrt(c/2).
% At P = (1.10, 0.55): grad f = (2x, 4y) = (2.20, 2.20), i.e. the 45-degree direction,
% which is indeed normal to the ellipse there (the ellipse tangent has slope -x/(2y) = -1).
\begin{tikzpicture}[>=Latex, scale=1.5]
  \draw[->] (-2.35,0) -- (2.5,0) node[right, font=\footnotesize] {$x$};
  \draw[->] (0,-1.75) -- (0,1.95) node[above, font=\footnotesize] {$y$};

  % Three level sets of x^2 + 2y^2. The middle one, c = 2.142, is chosen to pass through the
  % marked point P = (1.10, 0.6826): 1.10^2 + 2*0.6826^2 = 1.21 + 0.932 = 2.142.
  \foreach \c in {0.9, 2.142, 4.0}{
    \draw[gray!70, thin] (0,0) ellipse ({sqrt(\c)} and {sqrt(\c/2)});
  }
  % Kept in the upper LEFT: the gradient arrow and its label occupy the upper right.
  \node[gray!75!black, font=\scriptsize, anchor=west] at (-2.25,1.45) {level sets $\{f=c\}$};

  % a path crossing the level sets
  \draw[blue!75!black, thick, domain=0.15:2.05, samples=60, smooth]
    plot (\x, {0.28 + 0.30*\x + 0.06*\x*\x});
  \node[blue!75!black, font=\scriptsize, anchor=south east] at (0.55,0.42) {$\gamma$};

  % the point P on the middle ellipse (c = 1.815) and on the path:
  % path at x = 1.10 gives y = 0.28 + 0.33 + 0.0726 = 0.6826 -- adjust c to match the point.
  \coordinate (P) at (1.10,0.6826);
  \fill[black] (P) circle (0.8pt);
  \node[font=\scriptsize, below right] at (1.13,0.63) {$\gamma(t)$};

  % gradient of f at P: (2x, 4y) = (2.20, 2.73), normalised and scaled to length 0.62
  \draw[red!80!black, very thick, ->] (P) -- ++(0.39,0.48)
    node[anchor=south west, font=\scriptsize] {$\nabla f$};
  % velocity: gamma'(x) = 0.30 + 0.12x = 0.432, direction (1, 0.432), scaled
  \draw[blue!75!black, very thick, ->] (P) -- ++(0.57,0.246)
    node[anchor=north west, font=\scriptsize] {$\gamma'$};

  \node[font=\scriptsize, align=center, anchor=north] at (0,-1.85)
    {$\dfrac{d}{dt}f(\gamma(t)) = \langle \nabla f, \gamma'\rangle$: only the component of
     $\gamma'$ along $\nabla f$ contributes,\\
     so $f$ is constant along a level set and grows fastest along $\nabla f$};
\end{tikzpicture}
\end{center}

% Sonnet 5 (Medium)


% --- exercises relocated here from the dissolved sheet 4 ---

% Originally: content/week-04.tex, end of the chapter

\begin{aiexercise}[Directional Derivative along a Unit Vector]
\label{ex:ai_directional_derivative}
% Generator: Gemini 3.6 Flash (Medium)
Let $f(x,y) := x^2 y$ defined on $\mathbb{R}^2$. Compute the directional derivative $D_v f(1,2)$ of $f$ at the point $(1,2)$ along the unit vector $v := \frac{1}{\sqrt{2}}(1,1)\transp$.
\end{aiexercise}

% Originally: content/week-04.tex, \exercisesheet{4}
% Source: exercises/Ex4_Analysis2_eng.pdf

\begin{exercise}[4.6 --- A directional derivative vanishes \textnormal{(important)}]
\label{ex:4.6}
Let $u \in C^1(\mathbb{R}^n)$ and $\nu \in \mathbb{R}^n$. Show that
\begin{enumerate}[label=\textbf{\arabic*.}]
  \item If $\partial_1 u \equiv 0$ then ``$u$ does not depend on $x_1$''; more rigorously: there
    exists a unique function $v \in C^1(\mathbb{R}^{n-1})$ such that
    \begin{equation*}
    u(x_1,\dots,x_n) = v(x_2,\dots,x_n) \quad \text{for all } x \in \mathbb{R}^n. \tag{2}
    \end{equation*}
  \item If $\partial_\nu u \equiv 0$ and $\nu\cdot e_1 \neq 0$ then ``$u$ is a function of $n-1$
    variables''; more rigorously: there exists a unique function $w \in C^1(\mathbb{R}^{n-1})$
    such that
    \[u(x_1,\dots,x_n) = w\!\left(x_2 - \frac{x_1\nu_2}{\nu_1}, \dots, x_n - \frac{x_1\nu_n}{\nu_1}\right) \quad \text{for all } x \in \mathbb{R}^n.\]
  \item \textbf{(*)} What can we conclude if we assume only that $\partial_1 u = 0$ in an open
    connected subset $U \subset \mathbb{R}^n$?
\end{enumerate}
\textit{Official hints:} \textbf{4.6.2} --- apply the mean value theorem to $u(x + t\nu)$.
\textbf{4.6.3} --- consider the domain $U := \{(x,y) : x^2 < y < x^2+1\}$ and the function
\[u(x,y) := \begin{cases} \max\{0, y-2\}^2 & \text{if } x \geq 0, \\ 0 & \text{if } x \leq 0.\end{cases}\]
\end{exercise}

% Originally: content/week-04.tex, \section{Solutions}

\section{Solutions}
\label{sec:chain_rule_solutions}

\begin{exercisesolution}[Solution to \cref{ex:4.6}]
% Generator: Claude Sonnet 5 (Medium)
\begin{enumerate}[label=\textbf{\arabic*.}]
  \item Define $v(x_2,\dots,x_n) := u(0,x_2,\dots,x_n)$, which is $C^1$ as the restriction of $u$
    to the hyperplane $x_1 = 0$. For fixed $(x_2,\dots,x_n)$, the function
    $t \mapsto u(t,x_2,\dots,x_n)$ has derivative $\partial_1 u \equiv 0$ everywhere, hence is
    constant in $t$; evaluating at $t=0$ gives $u(x_1,\dots,x_n) = u(0,x_2,\dots,x_n) =
    v(x_2,\dots,x_n)$ for all $x_1$. Uniqueness: if some $\tilde v$ also satisfies
    $u(x) = \tilde v(x_2,\dots,x_n)$, then setting $x_1 = 0$ gives $\tilde v = v$.
  \item Since $\nu\cdot e_1 = \nu_1 \neq 0$, consider the invertible linear change of coordinates
    $y_i := x_i - \dfrac{x_1\nu_i}{\nu_1}$ for $i = 2,\dots,n$, together with $x_1$ itself.
    Moving $x_1$ by $t$ while holding $y_2,\dots,y_n$ fixed forces $x_i$ to move by
    $t\nu_i/\nu_1$, i.e.\ moves $x$ in the direction $\tfrac{1}{\nu_1}\nu$. By the hint, applying
    the mean value theorem to $u(x+t\nu)$ shows $t \mapsto u(x+t\nu)$ is constant since
    $\partial_\nu u \equiv 0$; consequently $u$, expressed in the coordinates $(x_1, y_2,\dots,
    y_n)$, has vanishing $x_1$-derivative and is therefore (by part 1) independent of $x_1$ in
    these coordinates. Hence, there is a (unique, by the same argument as in part 1) function
    $w \in C^1(\mathbb{R}^{n-1})$ with
    \[
    u(x_1,\dots,x_n) = w\!\left(x_2 - \frac{x_1\nu_2}{\nu_1},\dots,x_n - \frac{x_1\nu_n}{\nu_1}\right).
    \]
  \item In general, only \textbf{locally}: on any open subset where every fiber
    $\{x_1 : (x_1,x_2,\dots,x_n) \in U\}$ (for fixed $(x_2,\dots,x_n)$) is connected, the same
    argument as in part 1 shows $u$ is locally of the form $v(x_2,\dots,x_n)$. But this need not
    hold \textbf{globally} on all of $U$, even though $U$ itself is connected: with
    $U := \{(x,y) : x^2 < y < x^2+1\}$ and $u(x,y) := \max\{0,y-2\}^2$ for $x \geq 0$, $u := 0$
    for $x \leq 0$, one checks $\partial_1 u \equiv 0$ throughout $U$ (near $x=0$ both branches
    agree, since $y < 2$ there), yet $u$ is not globally independent of $x$: for $y > 2$, the
    fiber of $U$ over $y$ splits into two disconnected intervals (one with $x>0$, one with
    $x<0$), and $u$ takes different values on each. So, connectedness of $U$ alone does not
    suffice; what is needed is connectedness of each vertical fiber of $U$.
\end{enumerate}
\end{exercisesolution}

\begin{exercisesolution}[Proof of \cref{ex:ai_directional_derivative}]
% Generator: Gemini 3.6 Flash (Medium)
The gradient of $f(x,y)$ is $\nabla f(x,y) = (2xy, x^2)\transp$. Evaluating at $(1,2)$ gives $\nabla f(1,2) = (4,1)\transp$. Since $f$ is continuously differentiable ($C^\infty$), the directional derivative along the unit vector $v$ is given by the inner product:
\[
D_v f(1,2) = \nabla f(1,2) \cdot v = \begin{pmatrix} 4 \\ 1 \end{pmatrix} \cdot \frac{1}{\sqrt{2}}\begin{pmatrix} 1 \\ 1 \end{pmatrix} = \frac{4 + 1}{\sqrt{2}} = \frac{5}{\sqrt{2}}.
\]
\end{exercisesolution}

